\documentclass[11pt]{article}
\usepackage{fontspec}
\setmainfont{Times New Roman}
\setsansfont{Arial}
\setmonofont{Consolas}
\usepackage{amsmath,amssymb}
\usepackage{geometry}
\usepackage{microtype}
\geometry{a4paper,margin=3cm}
\setlength{\parindent}{1.25em}
\setlength{\parskip}{0pt}
\emergencystretch=30em
\allowdisplaybreaks
\raggedbottom
\newcommand{\ideal}[1]{\mathfrak{#1}}

\begin{document}

% Unit: Paper 19 title, contents, and introduction. Deterministic Interslavic Cyrillic reader generated from the Latin source and manually checked.

\setcounter{footnote}{0}
\section*{19. Теорија идеалов в колцевых областјах}
\begin{center}
Math. Ann. 83 (1921), стр. 24--66
\end{center}

\begin{center}\textbf{Содржанье.}\end{center}
\noindent Увод.\\
\S~1. Колцева област, идеал, условје конечности.\\
\S~2. Представјенье идеала како најменшего обчего многократного конечно многых неразложимых идеалов.\\
\S~3. Равност числа компонентов при двох различных разложеньах на неразложиме идеалы.\\
\S~4. Примарне идеалы. Јединост придружених простых идеалов при двох различных разложеньах на неразложиме идеалы.\\
\S~5. Представјенье идеала како најменшего обчего многократного највечих примарных идеалов. Јединост придружених простых идеалов.\\
\S~6. Једино представјенье идеала како најменшего обчего многократного релативно-просто неразложимых идеалов.\\
\S~7. Јединост изолованых идеалов.\\
\S~8. Једино представјенье идеала како продукта взаимно-просто неразложимых идеалов.\\
\S~9. Разширјенье изследованьа на модулы. Равност числа компонентов при разложеньах на неразложиме модулы.\\
\S~10. Специалны случај области полиномов.\\
\S~11. Примеры из теорије чисел и из теорије диференциалных изразов.\\
\S~12. Пример из теорије елементарных делительев.

\subsection*{Увод.}
Содржанье настојечеј работы јест \emph{пренос разложных теоремов целых рационалных чисел, односно идеалов в алгебраичных числовых польах, на идеалы в любих областјах целости, обчејше в колцевых областјах}. За разуменье того преноса најпрво наведемо разложне теоремы за целе рационалне числа в форме неколико различнеј од обычној формулације.

Ако в
\[
 a=p_1^{\varrho_1}p_2^{\varrho_2}\cdots p_\sigma^{\varrho_\sigma}=q_1q_2\cdots q_\sigma
\]
примарне степени \(q_i\) берут се како компоненты разложеньа, тогда ты компоненты имајут следујуче характеристичне својства:

1. Оне сут \emph{попарно взаимно просте}; але ниједно \(q_i\) не може быти представјено како продукт попарно взаимно простых чисел, значи в том смыслу имаје место неразложимост. Из попарној взаимној простости далье следује, же продукт \(q_1\cdots q_\sigma\) става равны најменшему обчему многократному \([q_1\cdots q_\sigma]\).

2. Кажде двије компоненты, \(q_i\) и \(q_k\), сут \emph{релативно просте}; т. ј. ако \(bq_i\) јест делимы чрез \(q_k\), тогда \(b\) јест делимы чрез \(q_k\). И в том смыслу имаје место неразложимост.

3. Каждо \(q\) јест \emph{примарно}; т. ј. ако продукт \(b\cdot c\) јест делимы чрез \(q\), але \(b\) не јест делимы, тогда некоја степень\footnote{Ако та степень всегда јест прва, тогда, како знано, јде о просте числа.} од \(c\) јест делимаја. Представјенье јест далье тако представјенье чрез \emph{највечше примарне компоненты}, понеже продукт двох различных \(q\) уже не јест примарны. Такоже в односу к разложеньу на највечше примарне компоненты та \(q\) сут неразложиме.

4. Каждо \(q\) јест \emph{неразложимо} в том смыслу, же јего не можно представити како најменше обче многократно двох властных делительев.

Свез тых примарных чисел \(q\) с простыми числами \(p\) состоји в том, же к каждому \(q\) сушчствује једно и, кроме знака, само једно \(p\), кторо јест делитель \(q\) и чија степень јест делимаја чрез \(q\): придружено просто число. Ако \(p^\varrho\) јест најнижша така степень -- \(\varrho\) експонент од \(q\) --, тогда ту особливо \(p^\varrho\) јест равно \(q\). Теорем јединости можно ныне изректи тако:

\emph{При двох различных разложеньах целого рационалного числа на неразложиме, највечше примарне компоненты \(q\) совпадајут число компонентов, придружене просте числа (до знака) и експоненты. Понеже \(p^\varrho=q\), из того следује такоже совпаданье самых \(q\) (до знака).}

Неопредельеност дана знаком, како знано, одстраньа се, ако вместо чисел разглядају се изведене из ньих идеалы (все числа делимы чрез \(a\)); тогда формулација важи точно тако за једино разложенье идеалов конечных алгебраичных числовых поль на степени простых идеалов.

В дальнем (\S~1) за основу се бере обча колцева област, ктора мора изполньати само условје конечности: кажды идеал области имаје конечны идеалны базис. Без такого условја конечности неразложиме и просте идеалы не мусили бы обче сушчствовати, како показује област всех целых алгебраичных чисел, в которој не имаје разложеньа на просте идеалы.

Показује се, же -- одповедајуче четырим характеристичным својствам компонентов \(q\) -- в обчем случају сушчствујут четыри одвојене разложеньа, кторе једно из другого настајут чрез подраздельенье. При том разложенье на взаимно-просто неразложиме идеалы јест продуктно представјенье, а остале три разложеньа сут редуциране (\S~2) представјеньа како најменша обча многократна. Свез меджу примарным идеалом -- и неразложиме идеалы сут примарне -- и придруженым простым идеалом такоже остаје: к каждому примарному идеалу \(\ideal Q\) једнозначно јест определьен придружены просты идеал \(\ideal P\), кторы јест делитель \(\ideal Q\) и чија степень става делимаја чрез \(\ideal Q\). Ако \(\ideal P^\varrho\) јест најнижша така степень -- \(\varrho\) експонент од \(\ideal Q\) --, тогда ту \(\ideal P^\varrho\) не треба совпадати с \(\ideal Q\). Теорем јединости изговарја се тако:

Разложеньа 1 и 2 сут једине; при двох различных разложеньах 3 или 4 совпадајут число компонентов и придружене просте идеалы;\footnote{Верјетно понад то совпадајут такоже експоненты, и јеште обчејше одповедајуче компоненты сут изоморфне.} изоловане идеалы (\S~7), кторе наступајут меджу компонентами, сут једнозначно определьене.

За доказ разложных теоремов из условја конечности выводит се најпрво ``теорем о конечној цепи'', перво изречены Дедекиндом за конечне числове модулы, и из ньега представјенье 4 каждого идеала како најменшего обчего многократного конечно многых неразложимых идеалов. Преобликованьем појма разложимости компонента из того се добива основны теорем јединости за разложенье 4 на неразложиме идеалы. Собранием по конечно многых компонентов се потом приходит к осталым разложеньам, чији теоремы јединости следут како последок из теорема јединости 4.

Напоследку показује се (\S~9), же представјенье чрез конечно много неразложимых компонентов имаје место и под слабшими предпоставками; комутативност колцевој области не јест потребна, и достатно јест вместо идеала разглядати модул в односу к области. В том обчејшем случају јеште важи равност числа компонентов при двох различных разложеньах, меджутым појмы просты и примарны сут связане с комутативностју и с појмом идеала; напротив, појем взаимној простости при идеалах в некомутативных областјах остаје.

Најпростејша колцева област, в которој фактично наступајут четыри одвојене разложеньа, јест област всех полиномов од \(n\) променных с льубыми комплексными коефициентами. Једина разложеньа можно ту ирационално толковати чрез поведенье алгебраичных творб, и теорем јединости за придружене просте идеалы одповедаје фундаменталному теорему теорије елиминације о јединој разложимости алгебраичных творб на неразложиме. Дальнье примеры дајут все конечне области целости из полиномов (\S~10). Але и проста област всех парных чисел, обчејше всех чисел делимых чрез једно устальено число, уже дава пример честично одвојеных разложень (\S~11). Пример теорије идеалов в некомутативных областјах дава теорија елементарных делительев (\S~12), гдје сушчствује једино разложенье на неразложиме идеалы, односно класы. Те неразложиме класы полно карактеризујут неразложиме составне чести елементарных делительев и можно в областјах, гдје обычна теорија елементарных делительев не успева, сматрјати јих како јеј еквивалент.

О сушчствујучеј литератури треба приметити следујуче: разложенье на највечше примарне идеалы за област полиномов с льубыми комплексными, односно целыми, коефициентами дал Lasker, и Macaulay јего в појединых точках далье развил.\footnote{E. Lasker, Zur Theorie der Moduln und Ideale. Math. Ann. 60 (1905), стр. 20, Satz VII und XIII. -- F. S. Macaulay, On the Resolution of a given Modular System into Primary Systems including some Properties of Hilbert Numbers. Math. Ann. 74 (1913), стр. 66.} Оба опирајут се на теорију елиминације, значи користајут факт, же полином можно једнозначно представити како продукт неразложимых полиномов. В самој ствари разложне теоремы за идеалы сут независне од тој предпоставкы, како може се очекивати из теорије идеалов в алгебраичных числовых польах и како показује настојеча работа. И примарны идеал при Laskeru и Macaulayu определьен јест на основе појмов из теорије елиминације.

Разложенье на неразложиме идеалы и разложенье на релативно-просто неразложиме здаје се, же в литератури не было замечено ни за област полиномов; само при Macaulayu находи се приметка о јединости изолованых примарных идеалов.

Разложенье на взаимно-просто неразложиме идеалы за област полиномов дал Schmeidler,\footnote{W. Schmeidler, Über Moduln und Gruppen hyperkomplexer Größen. Math. Zeitschr. 3 (1919), стр. 29.} користечи теорију елиминације за доказ конечности. Але там теорем јединости изречен јест само за класы идеалов, не за саме идеалы. Последньи теорем јединости находи се в обчеј работе,\footnote{E. Noether--W. Schmeidler, Moduln in nichtkommutativen Bereichen, insbesondere aus Differential- und Differenzenausdrücken. Math. Zeitschr. 8 (1920), стр. 1.} гдје јде о идеалы в некомутативных полиномовых областјах. Там се употребја само конечны идеалны базис; теоремы и методы зато оставајут важне за обче колцеве области, и настојеча работа јих заострује в односу к равности числа (\S~11). Настојеча изследованьа представльајут силно обобченье и дальнье развијанье појмовых творб, лежечих в основе обех работ. Сушчствено в обех работах јест преход од представјеньа како најменшего обчего многократного к адитивному разложеньу система класов остатков. Ту, дльа простејшего изложеньа, зново остајемо при најменшем обчем многократном; адитивному разложеньу тогда одповедаје преобликованье појма разложимости в својство комплемента (\S~3). Але по размыслам, кторе сушчствено одповедајут тым из обчеј работы, все наведене теоремы можно такоже разумети како адитивне разложне теоремы за систем класов остатков и некоторе подсистемы. Тај систем класов остатков твори колцо такој самој обчности како првоначално основано; именно каждо колцо можно разумети како систем класов остатков того идеала, кторы одповедаје целости идентичных релациј меджу елементами колца; или такоже некој подсистемы тых релациј, ако остале релације пријимајут се како уже изполньене в области.

Та приметка дава такоже место работам Fraenkela.\footnote{A. Fraenkel, Über die Teiler der Null und die Zerlegung von Ringen. J. f. M. 145 (1914), стр. 139. Über gewisse Teilbereiche und Erweiterungen von Ringen. Habilitationsschrift, Leipzig, Teubner, 1916. Über einfache Erweiterungen zerlegbarer Ringe. J. f. M. 151 (1920), стр. 121.} Fraenkel разгляда адитивне разложеньа колц, кторе подвргајут се таким ограничным условјам (сушчствованье регуларных елементов, дельенье чрез ньих, условје разложимости), же за одповедајучи идеал четыри разложеньа совпадајут. Из-за того совпаданьа јего условје конечности, же идеал има само конечно много властных делительев -- од чего се в чести одступаје -- не значи силнејше ограниченье неж наше. Изходишче Fraenkela јест условјено иными, в сучности алгебраичными, цельами јего работ; чрез алгебраично разширјенье он потом приходит к обчејшим колцам с мене ограничными условјами.

\clearpage

\end{document}

